% Direct-authority diplomatic French for printed p.215.
\clearpage
\SourcePage{215}

\LayerTarget{tohoku:comparison:equivariant-hom-stalk-to-fiber}
\LayerTarget{tohoku:equation:5.6.9}
\noindent bijectifs lorsque ${\mathbf O}$ est un faisceau \textit{cohérent} d'anneaux
\textit{noethériens} à gauche, et $A$ est cohérent en tant que ${\mathbf O}$-module ; et
que sous ces conditions sur $X,G,{\mathbf O}$, le $U_x$-module $B(x)$ est injectif chaque
fois que $B$ est un objet injectif de ${\mathbf C}^{{\mathbf O}(G)}$. (On remplacera dans la
démonstration, les ${\mathbf O}$-Modules libres de 4.1. par des $G$-${\mathbf O}$-modules du type
$L(U)$ introduits dans la démonstration de prop. 5.1.1). On en conclut
immédiatement que \textit{les homomorphismes naturels}

\LayerTarget{tohoku:comparison:equivariant-ext-stalk-to-fiber}
\LayerTarget{tohoku:equation:5.6.10}
\[
\tag{5.6.10}
\mathbf{Ext}_{{\mathbf O},G}^n(A,B)(y)
\to \operatorname{Ext}_{U_x}^n(A(x),B(x))
\]
\textit{déduits des homomorphismes} (5.6.9) \textit{sont des
isomorphismes si on suppose le groupe d'homéomorphismes $G$ discontinu au sens
de \textup{5.3}, ${\mathbf O}$ un faisceau cohérent d'anneaux noethériens à gauche, et $A$ un
$G$-${\mathbf O}$-Module, cohérent en tant que ${\mathbf O}$-Module}. Cet énoncé contient à la fois
les théorèmes 4.2.2 et 5.3.1 (le faisceau d'anneaux constant $\mathbf Z$ étant
cohérent et noethérien !) et précise dans les cas les plus importants la
structure des $\mathbf{Ext}_{{\mathbf O},G}^n(A,B)$.

\LayerTarget{tohoku:spectral-sequence:equivariant-ext-composite-functors}
Des formules (5.6.4) et (5.6.5), on peut tirer autant de suites spectrales par
application du théorème 2.4.1. Nous ne nous intéresserons qu'à celles déduites
de (5.6.5), aboutissant à $\operatorname{Ext}_{{\mathbf O},G}^{*}(X;A,B)$. Il faut encore
vérifier que dans chacune de ces trois formules donnant $h_{{\mathbf O},G,A}$ comme un
foncteur composé, la condition d'acyclicité usuelle est vérifiée, cela résulte
du

\LayerTarget{tohoku:lemma:5.6.1}
\textsc{Lemme 5.6.1.}\quad \textit{Supposons que $B$ soit un objet
injectif de ${\mathbf C}^{{\mathbf O}(G)}$, alors $\mathbf{Hom}_{{\mathbf O},G}(A,B)$ est un faisceau flasque
sur $Y$, $\operatorname{Hom}_{\mathrm O}(A,B)$ est un $G$-module
$\Gamma^G$-acyclique, et $\mathbf{Hom}_{\mathrm O}(A,B)$ est un $G$-faisceau
$\Gamma_X^G$-acyclique.}

\LayerTarget{tohoku:proof:lemma-5.6.1}
En vertu de proposition 5.1.2, on est ramené au cas où $B$ est le faisceau
produit défini par une famille $(B_x)_{x\in X}$ de $U_x$-modules injectifs
$B_x$. Il en résulte que le $G$-faisceau abélien
$\mathbf H=\mathbf{Hom}_{\mathbf O}(A,B)$ est le faisceau produit défini par la famille
de $G$-modules
$H_x=\operatorname{Hom}_{{\mathbf O}(x)}(A(x),B(x))$.
\LayerTarget{tohoku:construction:product-hom-sheaf}
Alors $\mathbf{Hom}_{{\mathbf O},G}(A,B)$ est le faisceau produit sur $Y$ défini par la
famille des groupes
\LayerTarget{tohoku:equation:product-stalk-invariants}
$H_y=\displaystyle\prod (H_x)^{G_x}$ (produit étendu aux $x$ tels que $f(x)=y$) donc est flasque.
Admettons un instant que les $H_x$ sont des $G_x$-modules
$\Gamma^{G_x}$-acycliques, (ce sera prouvé dans le corollaire ci-dessous).
\LayerTarget{tohoku:construction:coinduced-hom-module}
Alors le $G$-module $\overline H_x$ induit par le $G_x$-module $H_x$ par
extension contravariante des scalaires est $\Gamma^G$-acyclique comme bien
connu, donc
\LayerTarget{tohoku:equation:global-hom-product}
$\operatorname{Hom}_{\mathbf O}(A,B)=\Gamma_X\mathbf H
=\displaystyle\prod_{x\in X/G}\overline H_x
$ est $\Gamma^G$-acyclique. Enfin, $\mathbf H$ étant flasque donc $\Gamma_X$-acyclique, il résulte de la
deuxième suite spectrale du théorème 5.2.1 que
\LayerTarget{tohoku:equation:equivariant-hom-cohomology-vanishing}
$H^n(X;G,\mathbf H)=H^n(G,\Gamma_X\mathbf H)$, qui est nul pour $n>0$ d'après ce qu'on a vu. Cela achève la
démonstration du lemme 5.6.1, pourvu qu'on prouve le cas particulier suivant :

\LayerTarget{tohoku:corollary:coinduced-hom-acyclicity}
\textsc{Corollaire.}\quad \textit{Soit $O$ un anneau avec unité ou
opère un groupe $G$, soit $U$ l'anneau engendré par $O$ et $\mathbf Z(G)$,
soumis aux relations de commutation $g\lambda g^{-1}=\lambda^g$ pour $g\in G$,
$\lambda\in O$. Soient $A$ et $B$ deux $U$-modules, $B$ injectif. Alors le
$G$-module $\operatorname{Hom}_{\mathbf O}(A,B)$ est $\Gamma^G$-acyclique. (Ne pas croire
qu'il soit injectif !)}

\LayerTarget{tohoku:proof:coinduced-hom-acyclicity}
\LayerTarget{tohoku:definition:coinduced-module-b-prime}
Soit $B'=\operatorname{Hom}_{Z}(\mathbf Z(G),B)$ le groupe des
applications $f$ de $G$ dans $B$, nous

\TohokuContinuation{216}
